\section{Error Analysis and Discussion}

The main experimental observation is that the lexicon baseline is strongest on the current dataset. This is expected because many examples were designed to express emotions through explicit words such as \textit{happy}, \textit{furious}, \textit{afraid}, \textit{sad}, \textit{shocked}, or \textit{trust}. Such words are likely to have strong associations in the NRC Hashtag Emotion Lexicon.

The Naive Bayes model performs worse than the lexicon baseline. This is not surprising because it is trained from isolated lexicon terms, not from full manually labeled texts. The model can generalize through character n-grams, but its weak labels are noisy and its training data does not fully represent sentence-level context.

The hybrid model improves over Naive Bayes, showing that the learned signal can be useful when constrained by lexical evidence. However, it does not beat the pure lexicon method on this dataset. This suggests that, for a small evaluation set with explicit emotion words, the interpretable lexicon signal is more reliable than the weakly supervised learned signal. On a larger or noisier dataset, especially with misspellings, hashtags, and less direct vocabulary, the hybrid method may become more competitive.

The confusion matrices generated by the code can be used for more detailed inspection. Likely confusions include anticipation versus fear or surprise, and trust versus fear or sadness, because some words may have associations with multiple emotions depending on context.
